\scanpage{77}

C'est \unsure{associé|mené} à $B'(f,g)=\frac{dg}{g}-\frac{df}{f}$ .\,.

Si on veut par contre un invariant sur $\unsure{\mathscr E|\mathcal E}_X^*$, on prendra

\[
 g\longrightarrow\frac{g}{dg}\qquad \mathscr E_X\longrightarrow\underline\omega_{X/\mathbb C}^{\otimes(-1)}
\]

on \deleted{trouve} définit

\[
 \deleted{\unsure{B}}C^{!}(f,g)=\deleted{\frac{dg}{g}-\frac{df}{f}}\;g(dg)^{-1}-f(df)^{-1}
\]

satisfaisant à a) b) c) [$C^{?!}(f,g)=C^{?!}(f,ag)$ si $a\in\mathbb C^{\times}$) d) ($C^{?!}(f,g)=C^{?!}(f,h)$ et $g(x)=h(x)$ \uline{en un pt qui n'est pas un zéro} de $g$ ou $h$ $\Rightarrow g=h$ ] et e) \deleted{\unsure{$[g\mapsto C(f,g)$|$[g\mapsto C^!(f,g)]$}} pas tout à fait e, car $g(dg)^{-1}-f(df)^{-1}$ \unsure{p.ex|p.ex.} $g(dg)^{-1}$ \unsure{ne|va} pas être id. nul si $g\in\Gamma\,\mathscr E_X$ ! Donc il faut se borner à

\[
 \mathscr E_X^*=\underline{\mathrm{Ét}}_{\mathbb C\text{-an}}(X,\unsure{G_m^*|G_m^{\times}|\mathbb C^*})
\]

au quel cas on l'envoie dans $\underline\omega_{X/\mathbb C}^{\otimes(\deleted{\unsure{+}}-1)\,*}$ par $g(dg)^{-1}$, ou au choix dans $\underline\omega_{X/\mathbb C}^*$ par $dg/g$, et il y a lieu de considérer $\omega_{X/\mathbb C}^{\times}$ comme un torseur sous $\underline O_X^*$\unsure{,|;|.}

Donc ici on \unsure{poserait|prend}

\[
 C^{\unsure{?|\prime|\natural}}(f,g)=\frac{dg}{g}\mathbin{:}\frac{df}{f}=\frac{dg}{df}\,\frac{g}{f}\in\Gamma(X,\underline O_X^{\times})
\]

Dans tous ces cas, l'espace des valeurs \inserted{$(\mathbb P^1(\mathbb C),\mathbb E^1(\mathbb C),G_m(\mathbb C))$} pour les $f,g$ était un espace \uline{homogène} sous le groupe d'op. envisagé ($\unsure{\mathrm{Gb}(1)|\mathrm{GP}(1)|\mathrm{GL}(1)}$, $G_a$, $G_m$) ce qui semble condition de succès ($\unsure{\mathbb C^{\times}|G_m^{\times}|\mathbb C^*}$ opérant sur $\mathbb C$ : mauvais !) 

c/ Groupe $\mathrm{Aff}(1)$ \deleted{$\unsure{G}$} $=G_m\cdot G_a$ opérant sur $\mathbb E^1$.

On considère alors

% archive sheet number: 76
